\clearpage
\SourceLeaf{457}
\begin{center}\large EINFACHSTE ABLEITUNG DES\\GRUND-LEHRSATZES\end{center}
\begin{center}BETREFFEND\end{center}
\begin{center}\large DIE KÜRZESTEN LINIEN AUF REVOLUTIONSFLÄCHEN.\end{center}
\begin{center}\begin{tabular}{lp{.73\linewidth}}
$P,P'$ & zwei Punkte auf einer solchen Fläche\\
$r$ & Länge der kürzesten Linie zwischen $P$ und $P'$\\
$A,A'$ & Azimuthe derselben in $P$ und $P'$,\\
& letztere von Ost nach West gezählt\\
$\varphi,\lambda;\ \varphi',\lambda'$ & Breite und Länge der Punkte $P,P'$\\
$\rho,\rho'$ & Halbmesser der Parallelkreise durch $P$ und $P'$\\
$R,R'$ & Krümmungshalbmesser der Meridiane in diesen Punkten.
\end{tabular}\end{center}
\noindent $r+dr,\ \varphi+d\varphi,\ \lambda+d\lambda,\ \varphi'+d\varphi',\ \lambda'+d\lambda'$ Werthe, in welche $r,\varphi,\lambda,\varphi',\lambda'$ übergehen, indem man anstatt $P,P'$ zwei andere ihnen unendlich nahe Punkte substituirt.

Man hat dann
\[
dr=R\cos A.d\varphi-R'\cos A'.d\varphi'-\rho\sin A.d\lambda+\rho'\sin A'.d\lambda'.
\]
\noindent Offenbar wird aber $dr=0$, wenn
\[
d\varphi=0,\qquad d\varphi'=0,\qquad d\lambda'=d\lambda
\]
\noindent genommen wird. Also
\[
\rho\sin A=\rho'\sin A'.
\]
\begin{flushright}W. Z. B. W.\end{flushright}
\Needspace{7\baselineskip}
\begin{center}\textbf{BEMERKUNG.}\end{center}
\noindent\hspace*{1.2em}Die vorstehende Notiz, die sich in einem Handbuche befindet, stammt wahrscheinlich aus dem Anfange der vierziger Jahre.
\begin{flushright}\textsc{Stäckel}.\end{flushright}

